\section{Background}
\label{sec:background}

\subsection{Koopman operator theory}
\label{sec:koopman-background}

This subsection largely follows the review of \citet{brunton_modern_2022} with some changes in notation. \note{It's meant to be a bit more accessible; I haven't taken measure theory. Sorry.}

\paragraph{Koopman operator and observables.}
Consider a discrete-time nonlinear dynamical system
\begin{equation}
x_{t+1} = f(x_t), \qquad x_t \in \mathcal{X} \subset \mathbb{R}^{d_x},
\end{equation}
where $x_t$ is the state and $f : \mathcal{X} \to \mathcal{X}$ describes the dynamics.
Koopman operator theory studies the evolution of \emph{observables}, which are measurement functions $h : \mathcal{X} \to \mathbb{R}$.
The Koopman operator $\mathcal{K}$ is defined by composition:
\begin{equation}
(\mathcal{K}h)(x) = h(f(x)).
\end{equation}
Although $f$ may be nonlinear, $\mathcal{K}$ is always linear on the (typically infinite-dimensional) function space of observables. Then, instead of evolving the state $x_t$ directly, we can evolve a vector of observables
\begin{equation}
\phi(x) = \big[h_1(x), \dots, h_{d_z}(x)\big]^\top.
\end{equation}
If the span of $\{h_1,\dots,h_{d_z}\}$ is invariant under $\mathcal{K}$, 
then there exists a matrix
$K \in \mathbb{R}^{d_z \times d_z}$ such that
\begin{equation}
\phi(x_{t+1}) = \phi(f(x_t)) = \mathcal{K}\phi(x_t) = K\, \phi(x_t),
\end{equation}
so that the nonlinear dynamics of $x_t$ are exactly represented by a linear dynamical system
in the observable coordinates $z_t = \phi(x_t)$.
To learn linear latent-state models of nonlinear dynamics, the goal is to find observables that approximately span such a Koopman-invariant subspace.

\paragraph{Approximating the Koopman operator: DMD/eDMD.}
In practice, we do not have access to the infinite-dimensional Koopman operator, and we must work with a finite family of observables.
Given a feature map $\phi : \mathcal{X} \to \mathbb{R}^{d_z}$, one seeks a finite-dimensional approximation
\begin{equation}
\phi(x_{t+1}) \approx K\, \phi(x_t),
\end{equation}
where $K \in \mathbb{R}^{d_z \times d_z}$ is a learned linear model.
Data-driven approximations such as dynamic mode decomposition (DMD) \citep{schmid_dmd_2010, rowley_2009_spectral, tu_dynamic_2014} and extended DMD (eDMD) \citep{williams_datadriven_2015, williams_kernel-based_2016} estimate $K$ from pairs $\{(x_t, x_{t+1})\}$ by solving
\begin{equation}
K = \arg\min_{\hat K} \sum_{t} \big\| \phi(x_{t+1}) - \hat K\, \phi(x_t) \big\|_2^2,
\end{equation}
with $\phi$ chosen from a (possibly large) dictionary of basis functions.
If the chosen observables are approximately invariant under $\mathcal{K}$, then the learned $K$ captures a
nearly linear evolution of the feature vector $z_t = \phi(x_t)$ while still allowing rich nonlinear behaviour
in the original state $x_t$.

DMD is a special case: we take $\phi(x)=x$ itself, so DMD directly fits a linear map
\begin{equation}
K^{\mathrm{DMD}} = \arg\min_{\hat K} \sum_{t} \big\| x_{t+1} - \hat K x_t \big\|_2^2
\end{equation}
to the state.
This can be viewed as learning a representation in which the transition model is as close as possible to linear while still capturing important nonlinearity in the observed trajectories.


% \paragraph{Dynamical systems definitions.}
% Recall the key notions from nonlinear dynamics that will be important for our multi-basin problem.



% A \emph{non-wandering set} $S$ is a set of points with the property that for every neighborhood $U$ of any $x \in S$ and every $T > 0$, there exists a time $t \ge T$ and an initial condition $x_0 \in U$ such that the trajectory returns to $U$ at time $t$.
% Informally, trajectories starting in $S$ keep returning arbitrarily close to $S$.

% An \emph{attractor} $A$ is a non-wandering set that has the property of asymptotic stability: trajectories that start near the set approach it asymptotically as $t \to \infty$. The attractor has an open \emph{basin of attraction}
% \begin{equation}
% \mathcal{B}(A) = \big\{ x : \operatorname{dist}(\phi_t(x), A) \to 0 \text{ as } t \to \infty \big\},
% \end{equation}
% where $\phi_t(x)$ denotes the trajectory at time $t$ starting from $x$.
% The basin of attraction is just the set of initial conditions whose trajectories approach the attractor $A$ in the long-time limit.

\paragraph{Problems with multiple basins of attraction.}

While Koopman representations are powerful, they face theoretical limitations in systems with complex global topology. A key concept here is the \emph{basin of attraction}. First, an \emph{attractor} $A$ is a set of points that has the property of asymptotic stability: trajectories that start near the set approach it asymptotically as $t \to \infty$ 
(e.g. a \emph{fixed point} $x^\ast$ that satisfies $f(x^\ast) = x^\ast$). Then, the basin of attraction $\mathcal{B}(A)$ is the set of initial conditions whose trajectories approach the attractor $A$ in the long-time limit:
\begin{equation} 
\mathcal{B}(A) = \{ x \in \mathcal{X} : \lim_{t \to \infty} \text{dist}(f^t(x), A) = 0 \}.
\end{equation}

For many nonlinear systems, especially those with multiple basins of attraction, no finite-dimensional Koopman-invariant subspace exists that both (1) contains the original state coordinates $x$, and (2) is globally invariant under $\mathcal{K}$ in a way that yields exact dynamics on all basins simultaneously.

\citet{lan_linearization_2013}, also summarized in \citet{brunton_koopman_2016}, show that such finite-dimensional invariant subspaces containing the state can exist only for systems with a single isolated fixed point whose basin covers the state space.
% If a nonlinear system has more than one isolated fixed point, it cannot admit a finite-dimensional Koopman-invariant subspace that both spans the state and yields exact global dynamics.
% However, for each \emph{single} basin of attraction, we can linearize the dynamics within that basin nicely.
% Near a hyperbolic fixed point, classical linearization theory guarantees the existence of local coordinates in which the dynamics are well-approximated by a linear system, and Koopman eigenfunctions provide one way to construct such local coordinates.
% The only struggle is: how do we glue these separate basins, where we can linearize the dynamics separately, into one finite-dimensional global system that covers all basins and still keeps the dynamics linear and the map from latent $z$ to state $x$?
Any kind of finite-dimensional Koopman-learned latent for a multi-basin system must either:
\begin{itemize}
    \item be effectively \emph{piecewise linear}, where the latent includes some implicit ``which basin am I in?'' variable, or
    \item be \emph{approximate}, in the sense that it smooths over the boundaries between regions with different dynamics and is not exact anywhere.
\end{itemize}
This observation motivates our use of sparse coding to make explicit which local model is active, and periodic re-encoding to switch between locally valid linearizations as trajectories move across basin boundaries.



\paragraph{Koopman predictors for controlled systems (optional).}
For controlled systems
\begin{equation}
x_{t+1} = f(x_t, u_t),
\end{equation}
Koopman-based approaches often add controls to the higher-dimensional latent where the dynamics are linear:
\begin{equation}
z_{t+1} = A z_t + B u_t, \qquad x_t \approx C z_t,
\end{equation}
where $z_t = \psi(x_t)$ is a vector of learned observables, and $A,B,C$ are learned from data \citep{korda_linear_2018}.
This predictor has linear dynamics in the lifted space, but the feedback law $u_t = \kappa_{\mathrm{lift}}(z_t)$ induces a nonlinear controller $\kappa(x_t) = \kappa_{\mathrm{lift}}(\psi(x_t))$ in the original state.
Such linear predictors can be plugged into standard tools such as LQR and MPC while the nonlinearity is hidden in the fixed lifting map $x \mapsto z$ and decoder.

There are many benefits to having linear dynamics in the latent space.
Firstly, linear dynamics allow for LQR solvers to be applied to derive closed-form solutions to optimal control problems \citep{kalman_general_1960}.
Linear dynamics also significantly simplify system identification and interpretability, and can be efficiently unrolled with parallelism (e.g., in modern state-space models \citep{gu2022efficiently, gu_2023_mamba, smith2023simplified, dao_2024_mamba2}), which is attractive for long-horizon prediction and planning.



\subsection{Koopman autoencoders and periodic re-encoding}
\label{sec:background-reencoding}

Recent work \citep{lusch_deep_2018, azencot_forecasting_2020, shi_deep_2022, frion_leveraging_2023} has combined Koopman operator theory with deep autoencoders to learn nonlinear observables and linear latent dynamics directly from data, similar to DMD and eDMD.
In a typical Koopman autoencoder, an encoder $\phi : \R^{d_x} \to \R^{d_z}$ and decoder $\psi : \R^{d_z} \to \R^{d_x}$ are trained together with a linear latent Koopman transition $K$ so that
\begin{equation}
z_t \approx \phi(x_t), \qquad z_{t+1} \approx K z_t, \qquad x_t \approx \psi(z_t).
\end{equation}
We might do this by minimizing the following objective:
\begin{equation}
\min_{K, \phi, \psi} \sum_\mathcal{D} ||x_t - \psi(\phi(x_t)) ||_2 + \lambda \cdot || \phi(x_{t+1}) - K \phi(x_t)||_2
\end{equation} 
where $\lambda > 0$ is a scalar.
Once trained, one can simulate forward in time by iterating the linear latent dynamics and decoding:
\begin{equation}
z_{t+k} \approx K^k z_t, \qquad \hat x_{t+k} = \psi(z_{t+k}).
\end{equation}

\citet{fathi_course_2023} observed that unrolling the linear dynamics in latent space in this way leads to long-term drift in trajectories when mapped back to the state space.
Because the encoder and decoder are not exact inverses and the Koopman-invariance of the learned features is only approximate, latent trajectories $z_t$ can slowly drift away from the region where the local linear approximation is valid, resulting in trajectories that cross or develop spurious attractors when decoded to $x$.

To address this, they introduced \emph{periodic re-encoding}.
Rather than letting the latent state evolve indefinitely under $K$, one intermittently decodes and re-encodes:
\begin{equation}
z_{t+k} = \phi\!\left(\psi\!\left(K^k z_t\right)\right),
\end{equation}
for some re-encoding period $k$.
More concretely, the trajectory is evolved linearly in the latent space for several steps, decoded back to the original state, and then this decoded state is re-encoded to obtain a fresh latent representation.
This contrasts with a one-shot encoding at the initial time, where the same linear approximation is used regardless of how far the trajectory has moved in state space.

In the context of the result by \citet{lan_linearization_2013}, periodic re-encoding can be understood as a mechanism for switching between locally-valid linearizations as the dynamics traverse different regions or cross basin boundaries.
Within a single basin away from the boundaries, the encoder–decoder pair should behave approximately linearly along trajectories, but near basin boundaries, the geometry of invariant sets and eigenfunctions changes rapidly: observables that were approximately Koopman-invariant in one region may no longer span an invariant subspace in the new region.
By re-encoding at these times (or at a fixed period that is short compared to the time spent near boundaries), the model effectively refreshes the latent embedding so that the latent linear dynamics are again centered on the current part of the attractor or basin. This should help the system transition cleanly between different basins of attraction.

\subsection{Sparse coding}
\label{sec:background-sparse}

Sparse coding seeks a representation $z$ such that $x \approx Dz$ for a learned dictionary $D : \R^{d_z} \to \R^{d_x}$, via the optimization
\begin{equation}
\min_z \;\; \frac{1}{2} \|x - D z\|_2^2 + \lambda \|z\|_1,
\end{equation}
where $\lambda > 0$ controls the sparsity of $z$ \citep{olshausen_emergence_1996, chen_atomic_1998, donoho_optimally_2003}.
The $\ell_1$ penalty induces a mapping $x \mapsto z(x)$ that is piecewise linear and piecewise constant in its support: on each region of input space, the same small set of dictionary atoms is active, and within that region, $z(x)$ depends linearly on $x$.
Geometrically, this models the latent space as a union of lower-dimensional linear subspaces, with each subspace indexed by the support (active set) of the sparse code.

At a high level, this union-of-subspaces structure is exactly the kind of structure we would like to impose on a multi-basin nonlinear control problem.
Each support pattern corresponds to a locally linear regime; moving between regimes corresponds to changing which atoms are active.

In our setting, we want to treat the Koopman features as a sparse combination of dictionary atoms.
The dictionary $D$ is global, and the support pattern of the sparse code $z$ acts like a mode label indicating which basin or local region we are currently in.
Inside a region (basin of attraction) where the dynamics look locally linear, that region uses the same dictionary atoms and therefore has its own local Koopman-like dynamics.
Crossing a basin boundary is equivalent to changing the sparse code support to another region with a different local linear model. I felt that this was also explained nicely in Appendix D of \citet{fathi_course_2023}.

This gives a natural way to build an adjacency matrix for a graph that encapsulates the nonlinear problem: nodes correspond to local regions of constant support, and edges correspond to observed transitions between supports along trajectories.
On each node, one can learn a local Koopman linear operator and solve the associated LQR-style control problem in closed form.
Planning trajectories from one side of the field to the other then becomes a graph traversal problem over these sparse-code regions.

\paragraph{Sparse coding and dimensionality.}
Sparse coding is not inherently ``better'' in higher dimensions; it helps when high-dimensional data lie near a union of low-dimensional subspaces and the representation is sufficiently sparse.
In that regime, sparsity can mitigate the curse of dimensionality by reducing the effective degrees of freedom and improving interpretability.
However, higher ambient dimension also increases sample complexity and makes support selection harder (and in our setting, can dilute basin-norm signals that drive exclusivity), so sparsity is only advantageous if the data are truly compressible and the optimization/regularization are well tuned.

\paragraph{LISTA as a fast sparse encoder.}
Classical sparse coding solves the above optimization with an iterative algorithm (e.g., ISTA/FISTA, coordinate descent), which is often too slow to use inside a learned dynamical model.
LISTA (Learned ISTA) \citep{gregor_learning_2010} replaces the iterative solver with a trainable, finite-depth network that approximates the sparse solution.

In LISTA, one learns a non-linear encoder $z = f_\theta(x)$ whose architecture mimics a truncated iterative shrinkage/thresholding algorithm.
The key nonlinearity is a component-wise shrinkage function with learned thresholds; by construction, the mapping $x \mapsto z(x)$ is explicitly piecewise affine, with regions defined by which coordinates survive thresholding.
This builds a strong inductive bias for \emph{mode selection} into the encoder: for each $x$, only a few coordinates of $z$ are active, picking a small subset of dictionary atoms that explain the input.
Compared to a generic nonlinear encoder where all latent coordinates can participate everywhere, LISTA-style encoders encourage a sparse, structured partition of state space into locally linear regimes, which is exactly what we will leverage in our Koopman control architecture.


\begin{algorithm}[ht]
\caption{LISTA Encoder (Forward Pass)}
\label{alg:lista}
\begin{algorithmic}[1]
\REQUIRE $x \in \mathbb{R}^{d_x}$; depth $T \in \mathbb{N}$; encoder $\phi$ with learnable parameters $W_e \in \mathbb{R}^{d_z \times d_x}$, $S \in \mathbb{R}^{d_z \times d_z}$, $\theta \in \mathbb{R}^{d_z}_{\ge 0}$
\ENSURE Output $z \in \mathbb{R}^{d_z}$ ($z$ is latent code $\phi(x)$)
\STATE $z^{(0)} \gets h_\theta (W_e x)$
\FOR{$t = 1$ \TO $T$}
    \STATE $c^{(t)} \gets W_e x + Sz^{(t-1)}$ 
    \STATE $z^{(t)} \gets \operatorname{h}_{\theta} \big(c^{(t)}\big)$
\ENDFOR
\RETURN $z^{(T)} = \phi(x)$
\end{algorithmic}
\end{algorithm}

We represent the decoder as a linear dictionary $D \in \mathbb{R}^{d_x \times d_z}$, so that $\hat{x} = D z$. The encoder $\phi(x)$ is implemented as a LISTA network (Alg. \ref{alg:lista}) with parameters $W_e \in \mathbb{R}^{d_z \times d_x}, S \in \mathbb{R}^{d_z \times d_z}$, and thresholds $\theta \in \mathbb{R}^{d_z}_{\ge 0}$; here $W_e$ plays the role of a learned operator analogous to $D^\top$, while $D$ is the dictionary used in reconstruction.

\section{Method}
\label{sec:method}

We focus on the deployment setting that matters for sparse Koopman models: trajectories come from nonlinear systems with multiple basins of attraction, but basin counts and basin labels are unknown at training time and unavailable at deployment. The latent support should therefore emerge from the model rather than from supervised mode labels. Our main comparison asks which sparse Koopman recipe gives the best long-horizon forecasts under periodic re-encoding, and when those same latents also recover basin structure without label supervision.

\subsection{Problem Setup}
Given transitions from a discrete-time system
\begin{equation}
    x_{t+1} = F(x_t), \qquad x_t \in \mathbb{R}^{d_x},
\end{equation}
we learn an encoder $\phi$, decoder $\psi$, and latent linear transition $K$ such that
\begin{equation}
    z_t = \phi(x_t), \qquad z_{t+1} \approx K z_t, \qquad x_t \approx \psi(z_t).
\end{equation}
The support pattern of $z_t$ is treated as a candidate regime descriptor. In multi-basin systems we do not expect a single finite-dimensional latent linearization to be globally exact, so the method is evaluated in the periodic-reencoding regime from Section~\ref{sec:background-reencoding}, not as an unconstrained open-loop rollout claim.

\subsection{Model Families}
Our main comparison centers on three model families.

\paragraph{Generic sparse MLP baseline.}
The main baseline is a Koopman autoencoder with generic MLP encoder and decoder and an $\ell_1$ penalty on the latent state. This model provides a strong forecasting anchor with minimal structural assumptions beyond latent sparsity.

\paragraph{Dense LISTA encoder.}
The first sparse-structured alternative replaces the MLP encoder with LISTA (Algorithm~\ref{alg:lista}), optionally preceded by a shallow nonlinear pre-transform, while keeping a dense learned Koopman matrix. The LISTA encoder preserves the piecewise-affine, support-selective inductive bias of sparse coding, but allows full latent coupling in the dynamics. This is the main cross-system LISTA comparator in the paper.

\paragraph{Block-diagonal LISTA.}
The second sparse-structured alternative keeps the same LISTA encoder and linear decoder but constrains the latent dynamics to a fixed block-diagonal operator,
\begin{equation}
    K_{\mathrm{block}} = \operatorname{blkdiag}\!\big(K^{(1)}, \dots, K^{(M)}\big).
\end{equation}
These blocks are model partitions, not ground-truth basins. We do \emph{not} assume one known basin per block or a known basin count at training time. The role of the block structure is only to reduce cross-coupling in the latent dynamics and test whether a more local linearization helps on hard oscillatory systems.

Earlier development also explored more aggressive structured Koopman variants and explicit regime-selection penalties. Those experiments were useful diagnostically, but they are not part of the main comparison because they did not produce a competitive or stable result. The final narrative below is driven by the generic sparse baseline, dense LISTA, and block-diagonal LISTA.

\subsection{Learning Objective}
\label{sec:learning-objective}

We train all main models from one-step pairs $(x_t, x_{t+1})$ without basin labels. The shared objective is
\begin{equation}
\label{eq:paper-main-loss}
\begin{aligned}
\mathcal{L} =
&\ \lambda_{\mathrm{rec}} \|x_t - \psi(\phi(x_t))\|_2^2
+ \lambda_{\mathrm{pred}} \|x_{t+1} - \psi(K\phi(x_t))\|_2^2 \\
&+ \lambda_{\mathrm{align}} \|\phi(x_{t+1}) - K\phi(x_t)\|_2^2
+ \lambda_{\mathrm{sp}} \Big(\|\phi(x_t)\|_1 + \|\phi(x_{t+1})\|_1\Big).
\end{aligned}
\end{equation}
The reconstruction term keeps the autoencoder on-manifold, the prediction and alignment terms encourage latent linear dynamics, and the sparsity term drives support selectivity. In targeted block-diagonal ablations we also tested block-usage regularizers, but the headline paper results do not depend on those auxiliary losses.

To keep the benchmark fair, the main comparisons use one frozen recipe per model family rather than system-by-system tuning. In particular, the dense LISTA result is produced by choosing a single external recipe on a small development subset and then rerunning the full benchmark without further per-system modification.

\subsection{Periodic Re-encoding and Forecast Metrics}

At inference time we evolve the latent dynamics for a fixed cadence $M$, decode, and re-encode:
\begin{equation}
    z_{t+M} \leftarrow \phi\!\left(\psi\!\left(K^M z_t\right)\right).
\end{equation}
This is the deployment mode of interest for multi-basin systems. It allows the latent representation to refresh when the learned linearization drifts or when the trajectory moves between local dynamical regimes.

Our primary forecasting metric is \texttt{H1000 best-periodic}: the system-median horizon-1000 error under the best saved periodic-reencoding cadence for that system. We also report pairwise system wins, good-system counts under the paper threshold \texttt{H1000 best-periodic} $< 10$, and fixed-cadence \texttt{periodic\_100} ablations to verify that the conclusions do not rely entirely on cadence selection. Autonomous every-step errors are treated as limitation diagnostics rather than the primary paper success criterion.

\subsection{Label-Free Basin Recovery Evaluation}
\label{sec:label-free-eval}

For benchmark systems with known basin labels, we test whether basin structure can be recovered from latent trajectories without using labels in training. For each checkpoint we encode 256 trajectories, construct per-trajectory features from both continuous latents and binary supports, reduce continuous views with PCA to 20 dimensions, and cluster trajectories with $k$-means using the oracle basin count only at evaluation time. We report Adjusted Rand Index (ARI) as the main clustering score. This protocol directly measures whether sparse support structure and latent geometry are basin-discriminative in a label-free setting.

\section{Experiments and Results}
\label{sec:experiments}

We organize the results around three questions. First, does dense LISTA become a credible benchmark competitor once training budget is matched fairly against the MLP anchor? Second, on the intrinsically hard oscillatory systems, is architecture the main lever or is step size the main lever? Third, when forecasting improves, do the latent supports also recover basin structure without supervision?

\subsection{Benchmark Protocol and Recipe Selection}

Our main benchmark contains 29 systems, uses latent size $d_z = 256$, sequence length $L=8$, and reports 3 seeds per system. The current paper claims are drawn from the 200k-step reruns, not from the earlier 50k matrix. The older matrix remains a useful symmetric audit, but it left a fairness hole because the MLP anchor had not yet been rerun at the longer budget.

We close that hole in two steps. First, we tune only the \emph{external} dense-LISTA recipe on a small 8-system development subset, keeping the dense architecture fixed. Second, we freeze the selected 200k dense recipe and rerun the full 29-system benchmark, while also rerunning the generic sparse MLP anchor at the same 200k budget. This turns the main dense-vs-MLP comparison into a matched-budget comparison rather than a 50k-vs-200k comparison.

Step size is handled by a fixed benchmark protocol rather than by ad hoc model rescue. When the default simulator step size is clearly too coarse, the benchmark halves $dt$ before spending more budget on architecture churn. This choice becomes central on Kuramoto and Hopfield.

\subsection{Main Benchmark: Dense LISTA Becomes Globally Competitive, but the MLP Retains the Best Median}

\begin{table}[ht]
\centering
\caption{Fair 200k benchmark summary on 29 systems. ``Systems won'' counts head-to-head wins against the fair 200k MLP anchor. Good systems satisfy \texttt{H1000 best-periodic} $< 10$.}
\label{tab:main-benchmark}
\begin{tabular}{lcccc}
\toprule
\textbf{Model family} & \textbf{Median H1000} & \textbf{Systems won} & \textbf{Good systems} & \textbf{Catastrophic} \\
\midrule
MLP sparse baseline & 0.0208 & --- & 25/29 & 3/29 \\
dense LISTA & 0.0232 & 18/29 & 26/29 & 2/29 \\
\bottomrule
\end{tabular}
\end{table}

Table~\ref{tab:main-benchmark} gives the central paper result. Once the budget is matched fairly, the generic sparse MLP baseline still has the best cross-system median \texttt{H1000 best-periodic} error ($0.0208$ versus $0.0232$ for dense LISTA). However, dense LISTA now wins a majority of systems head-to-head (18/29) and keeps one more system in the good-forecast band (26/29 versus 25/29). There are no systems where dense LISTA fails the good band while the fair 200k MLP anchor passes. Dense LISTA is therefore globally competitive, but not the best model by every summary metric.

This matched-budget result changes the interpretation of the earlier 50k matrix. In that symmetric 50k comparison, dense LISTA was close but still clearly second. After freezing one fair 200k recipe, its cross-system median improves from $0.0388$ to $0.0232$, which is strong evidence that much of the earlier dense-LISTA gap was optimization-limited rather than purely architectural. At the same time, the matched 200k rerun of the MLP anchor also improves substantially, so the comparison has to separate ``best overall median'' from ``wins more systems.'' Median alone would overstate the MLP advantage; win count alone would overstate the dense-LISTA advantage.

The cadence-selection question weakens, but does not erase, this benchmark story. Under a single global \texttt{periodic\_100} cadence, dense LISTA still wins 17/29 systems overall, but the good-system count falls to a tie at 22/29. The dense win-count advantage therefore survives fixed cadence, whereas the safety-margin advantage becomes weaker. For the final paper, this argues for reporting both the best-periodic headline and a fixed-cadence ablation, with full system-by-system tables deferred to the appendix.

Block-diagonal LISTA should not be framed as a global benchmark winner. Two matched 200k block-diagonal transfers reduce some isolated failures, but their cross-system medians remain much worse ($0.0555$ and $0.0477$) and they win only 7/29 and 5/29 systems against the fair MLP anchor. The block-diagonal result is therefore targeted, not global.

\subsection{Hard Systems: Smaller Step Size Is the Main Lever, and Block-Diagonal LISTA Helps Selectively}

The hard-system track should be read separately from the aggregate 29-system benchmark. On the intrinsic-HD systems, the main bottleneck is not simply encoder expressivity. It is the combination of intrinsic dimension, small-scale geometry, and numerical step size. The strongest empirical lever is reducing $dt$.

Kuramoto is the clearest positive case. In the dedicated 5-seed comparison at $N=16$, $dt=0.00625$, and 200k training steps, the MLP baseline reaches \texttt{H1000 best-periodic} $=27.02$, dense LISTA improves to $13.84$, and block-diagonal LISTA reaches $6.98$, cleanly entering the good-forecast regime. This ranking is not a cadence-selection artifact: under fixed \texttt{periodic\_100}, the same three-way ordering is reproduced exactly at horizon 1000.

\begin{table}[ht]
\centering
\caption{Kuramoto dimension sweep at $dt=0.00625$ and 200k steps. Entries are seed-median \texttt{H1000 best-periodic}.}
\label{tab:kuramoto-scaling}
\begin{tabular}{lccc}
\toprule
\textbf{Dimension $N$} & \textbf{MLP sparse} & \textbf{dense LISTA} & \textbf{block-diagonal LISTA} \\
\midrule
8 & 813.57 & 495.07 & 8.11 \\
16 & 30.18 & 13.44 & 7.07 \\
24 & 6.71 & 15.01 & 6.57 \\
32 & 6.68 & 92.28 & 5.92 \\
64 & 208.93 & 208.71 & 23.27 \\
\bottomrule
\end{tabular}
\end{table}

Table~\ref{tab:kuramoto-scaling} shows the limits of the hard-system claim. The smaller-$dt$ block-diagonal rescue is real through $N=32$, where block-diagonal LISTA is all-seeds-good at $N=16$, $24$, and $32$. But it is not fully robust at $N=8$, and it breaks by $N=64$, where the median rises to $23.27$ and only 2 of 5 seeds remain good. Dense LISTA does not transfer robustly to this regime: it improves over the MLP at $N=16$, but already falls behind by $N=24$ and is poor at $N=32$ and $64$.

The Kuramoto positive also survives a modest robustness shift. With a uniform frequency spread at $N=16$, $dt=0.00625$, and 200k steps, block-diagonal LISTA reaches \texttt{H1000 best-periodic} $=9.53$ while the MLP baseline stays at $44.46$. So the block-diagonal win is not a single-regime artifact. The limitation is elsewhere: even in these positive settings, autonomous every-step errors remain enormous, which means periodic re-encoding is still essential.

Hopfield tells a different story. Smaller $dt$ helps there too, but it does \emph{not} flip the architecture ordering. At $N=16$, $dt=0.00625$, and 200k steps, the MLP baseline reaches \texttt{H1000 best-periodic} $=3.36$ and block-diagonal LISTA reaches $8.82$. Both are inside the good band on the periodic-reencoding metric, but the MLP remains better, and every-step errors remain unstable for both. The hard-system conclusion is therefore narrow and clear: smaller $dt$ is the dominant lever, and block-diagonal LISTA is a selective Kuramoto rescue rather than a universal intrinsic-HD winner.

\subsection{Label-Free Basin Recovery Is Strong on Multiwell, Weak on Duffing, and Negative on Kuramoto}

The basin-support story is separate from the forecasting story. We therefore evaluate label-free basin recovery directly with the revised clustering protocol from Section~\ref{sec:label-free-eval}. This replaces the earlier trajectory-mean-only evaluation, which was too weak to reveal structure reliably.

\begin{table}[ht]
\centering
\caption{Summary of label-free basin recovery with the revised clustering protocol. Mean ARI is aggregated over feature views and seeds.}
\label{tab:label-free-summary}
\begin{tabular}{lccp{0.40\linewidth}}
\toprule
\textbf{System class} & \textbf{Basins} & \textbf{Mean ARI} & \textbf{Interpretation} \\
\midrule
Multiwell family (8 variants) & 5 & 0.794--0.991 & Strong positive: label-free clustering recovers basin structure across all variants, with near-perfect recovery on several systems. \\
Duffing & 2 & 0.19--0.24 & Weak positive: basin signal is present but sparse supports are not consistently aligned within basin. \\
Kuramoto & 5 & $\approx 0$ & Negative: latent supports and continuous features are not basin-separable under any tested view. \\
\bottomrule
\end{tabular}
\end{table}

Table~\ref{tab:label-free-summary} gives the clearest representational claim supported by the current evidence. Label-free basin recovery is real on the multiwell family: all 8 variants are strongly positive, with mean ARI between $0.794$ and $0.991$ and several seeds achieving essentially perfect recovery. The best ARI often comes from the MLP baseline, with LISTA families close behind, so the representational claim is broader than ``LISTA alone finds basins.'' Accordingly, we frame the representational result as follows: sparse Koopman autoencoders can recover basin structure without training-time basin labels on potential-well systems.

Duffing is the important weak case. The system still shows some label-free signal, but only weakly so. Mean ARI stays around $0.19$--$0.24$ across views and model families. The key point is that unique per-timestep supports do \emph{not} guarantee useful trajectory-level basin clustering. In the current Duffing checkpoints, within-basin support consistency is only about $12.8\%$ for one basin and $7.5\%$ for the other, so most trajectories do not share a stable basin-specific support even though the average representation is somewhat basin-discriminative.

Kuramoto is the strongest negative case and an important scientific caveat. Mean ARI is essentially zero across all views and model families. The supports are genuinely non-separable: within-basin and between-basin Hamming distances are almost identical (ratio $1.004$), the winding-number basin distribution is highly imbalanced, and nearly every trajectory has its own support pattern. A dedicated mode-support audit further confirms that this negative extends beyond clustering failure to the literal absence of reusable basin-specific supports. Although each winding-number basin technically has a ``unique'' modal support, this uniqueness is trivially degenerate: every trajectory activates a different support pattern, so modal supports are singletons that appear only once. Basin consistency is negligible, and threshold sweeps across multiple support definitions do not change the conclusion. The result is identical across all three model families. This is the clearest counterexample to any claim that good long-horizon forecasting automatically yields basin-support alignment. On Kuramoto, block-diagonal LISTA can win forecasting at small $dt$ while still failing completely as a basin-discriminative representation.

We do not report \texttt{competitive\_lv} here. The earlier setup was effectively one-basin, so those old support and clustering numbers are invalid for a multi-basin narrative, and the rebuilt multi-basin retrain was still incomplete.

\subsection{Summary of Supported Claims}

The current evidence supports a three-part narrative. First, fair optimization makes dense LISTA genuinely competitive on the main 29-system benchmark: it wins more systems and keeps slightly more systems in-band, but the MLP anchor still has the best overall median. Second, structured block-diagonal LISTA helps selectively on hard oscillatory systems, most clearly on smaller-$dt$ Kuramoto through $N=32$, but not as a universal benchmark solution and not through $N=64$. Third, label-free basin recovery is real but system-dependent: strong on multiwell, weak on Duffing, and negative on Kuramoto.

The main remaining limitation is autonomous stability. On the hard systems, every-step rollouts remain unstable even when periodic re-encoding forecasts are good. Accordingly, we treat forecasting under periodic re-encoding as the primary result, basin-support alignment as an empirical property that holds only on some system classes, and robust global linearization as outside the scope of the current evidence.

\appendix

\section{Legacy Exploratory Diagnostics}
\label{app:additional-diagnostics}

The appendix below records earlier structure-ablation and control-facing studies that were useful during development, but they are not part of the main narrative. The main claims of this draft are driven by the fair 200k benchmark, the smaller-$dt$ hard-system follow-ups, and the revised label-free clustering evaluation in the main text.

\subsection{Long-Horizon Dynamics: Spectral Radius as the Key Predictor}

We evaluated 25 checkpoints (5 target sizes $\times$ 5 Koopman structures including arrowhead without exclusivity) with long-horizon rollouts and eigenvalue analysis. The main result is consistent and strong: SR predicts long-horizon fate.

\begin{table}[ht]
\centering
\caption{Representative SR--stability outcomes (no re-encoding, horizon 1000).}
\label{tab:sr-stability}
\begin{tabular}{lccc}
\toprule
\textbf{Configuration} & \textbf{Max SR} & \textbf{H1000 MSE} & \textbf{Outcome} \\
\midrule
dense, ts=256 & 1.0262 & $6.78 \times 10^{17}$ & diverged \\
block-diagonal, ts=256 & 0.9996 & $3.58$ & stable \\
arrowhead (+excl), ts=128 & 0.9963 & $3.49$ & stable \\
arrowhead (no excl), ts=256 & 1.0215 & $1.52 \times 10^{13}$ & diverged \\
\bottomrule
\end{tabular}
\end{table}

Periodic re-encoding remains an effective fallback, typically recovering bounded long-horizon error ($\sim 3.6$--$3.9$ on Lyapunov) even when open-loop rollouts diverge.

\subsection{Sequence Training Does Not Stabilize Non-Diagonal Koopman Dynamics}

We tested two hypotheses under sequence training: (i) longer sequence length might stabilize $K$, and (ii) loss-weight tuning might recover stability.

\paragraph{Sequence-length sweep (48 runs).}
Across $L \in \{4,8,12,16\}$, target sizes $\{64,128,256\}$, and four Koopman structures, only diagonal $K$ is frequently stable, and even that degrades as $L$ increases. Dense, block-diagonal, and arrowhead are unstable in all tested settings.

\paragraph{Sequence-loss weight sweep (36 runs, block-diagonal, $L=8$, ts=128).}
No configuration achieved strict stability: \textbf{0/36} runs had SR $< 1$; SR ranged from 1.0014 to 1.2540. Setting prediction weight to 1.0 improved SR on average but did not cross the stability boundary. Therefore, weight tuning alone is insufficient; explicit spectral control is required.

\subsection{Objective Design for Basin Separability}

\paragraph{Single-loss ablation (24 runs).}
For block-diagonal $K$, usage-balance losses (\texttt{usage\_entropy}, \texttt{kl\_uniform}) consistently improve cosine separation over control, while strict one-block penalties (\texttt{low\_entropy}, \texttt{pairwise\_overlap}) often collapse.

\begin{table}[ht]
\centering
\caption{Best cosine separation by target size in block-loss ablation.}
\label{tab:block-loss-ablation}
\begin{tabular}{lccc}
\toprule
\textbf{Target size} & \textbf{Control sep.} & \textbf{Best loss (sep.)} & \textbf{Improvement} \\
\midrule
64 & 0.353 & \texttt{usage\_entropy} (0.786) & +0.433 \\
128 & 0.238 & \texttt{kl\_uniform} (0.852) & +0.614 \\
256 & 0.826 & \texttt{kl\_uniform} (0.884) & +0.058 \\
512 & 0.466 & \texttt{usage\_entropy} (0.791) & +0.325 \\
\bottomrule
\end{tabular}
\end{table}

\paragraph{Combined top-1 plus balance sweep (Phase 1, 72 runs, ts=256).}
The best seed-averaged configuration (\texttt{kl\_uniform}, top1 margin 0.05, one-block weight 0.3, balance weight 1.0) reaches cosine separation 0.8826, improving over the ts=256 control (0.8256). Across 36 two-seed configurations, 23/36 beat control separation, but sensitivity remains: 8/36 configurations lose full uniqueness in at least one seed, and the worst settings collapse sharply. Thus, combined objectives can help, but the robust hyperparameter region is narrow.

\subsection{Basin-to-Block Alignment vs. Basin-Unique Supports}

A key distinction from earlier iterations is that our primary objective is basin-unique sparse supports, not necessarily one-basin-per-block Koopman assignment. In this framing, results are strong: at target size 256, uniqueness is frequently saturated (13/13), and separation is high.

However, basin-to-block concentration from eigenvalue-correlation analyses remains highest at low capacity (e.g., concentration 0.87 at ts=64) and weak at higher capacity (around 0.10 at ts=256). This indicates that the model can achieve basin-discriminative sparse supports while distributing activity across multiple blocks. This behavior is compatible with our current objective, but it limits immediate extraction of one block per basin for controller assignment.

\subsection{LQR-Readiness and Architecture Selection}

The core control question is whether the learned latent regimes support \emph{reliable local LQR synthesis} under the intended label-free setting (unknown basin count/assignments at training and deployment). To test this directly, we ran a staged decision pipeline with pre-registered rules:
\begin{itemize}
    \item \textbf{Stage 1 (structure isolation):} compare a block-diagonal Koopman baseline against an arrowhead Koopman model \emph{without} exclusivity regularization on Lyapunov.
    \item \textbf{Stage 2 (practical head-to-head):} compare the selected block-diagonal baseline against a practical arrowhead model \emph{with} exclusivity and structured sparsity on Lyapunov.
    \item \textbf{Stage 3 (transfer):} repeat the Stage-2 comparison on Duffing.
\end{itemize}
All stages include basin-count misspecification through $B_{\text{proxy}}$ sweeps, so results do not rely on knowing the true number of basins at decision time.

For clarity, the \emph{DARE feasibility rate} is the fraction of discovered regimes for which the discrete algebraic Riccati equation admits a stabilizing solution (hence an LQR gain can be computed). The \emph{closed-loop stability rate} is the fraction of discovered regimes with spectral radius $\rho(A-BK_{\text{LQR}})<1$. The \emph{closed-loop cost reduction} compares finite-horizon quadratic cost under LQR versus open loop, reported as relative improvement.

\begin{table}[ht]
\centering
\caption{Staged LQR-readiness outcomes (label-free control metrics).}
\label{tab:lqr-readiness}
\begin{tabular}{p{0.24\linewidth}p{0.10\linewidth}c c c p{0.28\linewidth}}
\toprule
\textbf{Comparison} & \textbf{System} & \textbf{Runs} & \textbf{DARE feasibility} & \textbf{Closed-loop stability} & \textbf{Closed-loop cost reduction contrast} \\
\midrule
Stage 1: block-diagonal vs arrowhead (no exclusivity) & Lyapunov & 36 & 1.00 & 1.00 & Difference $= +0.0407$, 95\% CI $[-0.0888, 0.1650]$ (includes 0) \\
Stage 2: block-diagonal vs arrowhead (with exclusivity) & Lyapunov & 144 & 1.00 & 1.00 & Difference $= -11.5930$, 95\% CI $[-34.7496, 0.0006]$ (includes 0) \\
Stage 3: block-diagonal vs arrowhead (with exclusivity) & Duffing & 78 & 1.00 & 1.00 & Difference $= -0.0849$, 95\% CI $[-0.1100, -0.0632]$; magnitude $<0.10$ practical threshold \\
\bottomrule
\end{tabular}
\end{table}

These results address one key issue unambiguously: \emph{local-regime LQR is feasible and yields stable closed-loop dynamics in the discovered regimes}. We know this because both the DARE feasibility rate and the closed-loop stability rate are 1.0 across all compared arms in both systems, including misspecified $B_{\text{proxy}}$ settings.

At the same time, the architecture-selection issue is only partially resolved. Under the pre-registered decision rule (lexicographic ordering of DARE feasibility, then closed-loop stability, then closed-loop cost reduction, with a practical threshold of $+0.10$ absolute gain on feasibility or stability), no arm achieved a practical winning margin; therefore the final choice remained the simpler block-diagonal fallback.

Importantly, the arrowhead model shows a \emph{favorable trend} on the closed-loop control objective: in both Stage 2 (Lyapunov) and Stage 3 (Duffing), the cost-reduction contrast is in its favor, and in Stage 3 the confidence interval excludes zero. We therefore interpret arrowhead as the current cost-oriented contender, while block-diagonal remains the conservative default pending stronger, lower-variance confirmation.

Implication for the basin-switching control objective: the local-control component is now empirically supported, but the \emph{switching/gating} component remains open. The next experiments should therefore target closed-loop trajectories that cross regime boundaries, with periodic re-encoding used both as drift correction and as the mechanism that refreshes the latent regime assignment before controller switching.

\subsection{Exclusivity Ablation}

The arrowhead architecture pairs structured sparsity with an exclusivity regularizer that penalizes simultaneous activation of multiple basin blocks on a single input. Exclusivity was originally motivated as a mechanism to enforce one-basin-at-a-time encoding, thereby sharpening support--basin correspondence. However, the preceding LQR-readiness pipeline always coupled exclusivity with the arrowhead structure, confounding the two effects. We therefore ran a controlled ablation isolating the contribution of exclusivity alone.

The design holds all other hyperparameters fixed---structured layout, $\lambda_{\text{sparsity}} = 0.3$, $\lambda_{\text{global}} = 10^{-4}$, $\lambda_{\text{local}} = 10^{-3}$---and compares $\lambda_{\text{excl}} = 0.05$ (the value used in prior arrowhead runs) against $\lambda_{\text{excl}} = 0$ on the same Stage-2 Lyapunov grid: target size $\in \{128, 256, 512\}$, $B_{\text{proxy}} \in \{8, 13, 20\}$, eight seeds per cell (144 runs total, 129 completed).

\begin{table}[ht]
\centering
\caption{Exclusivity ablation: per-arm means across the Stage-2 Lyapunov grid.}
\label{tab:excl-ablation}
\begin{tabular}{l c c c}
\toprule
\textbf{Metric} & \textbf{With exclusivity} & \textbf{Without exclusivity} & \textbf{Effect} \\
\midrule
DARE feasibility (M2) & 1.000 & 1.000 & --- \\
CL stability (M3) & 1.000 & 1.000 & --- \\
CL cost reduction (M4) & $0.826 \pm 0.027$ & $0.823 \pm 0.028$ & $d = 0.07$, n.s. \\
Recovery success (M5) & $0.483 \pm 0.158$ & $0.549 \pm 0.159$ & $d = -0.82$, $p < 0.001$ \\
Cosine separation & $0.740 \pm 0.064$ & $0.801 \pm 0.036$ & $d = -0.97$, $p < 0.001$ \\
Intra-basin cosine & $0.990 \pm 0.005$ & $0.980 \pm 0.012$ & --- \\
Inter-basin cosine & $0.250 \pm 0.062$ & $0.179 \pm 0.031$ & --- \\
Max spectral radius (median) & 0.998 & 1.002 & --- \\
Fraction SR $> 1$ & 45\% & 61\% & --- \\
\bottomrule
\end{tabular}
\end{table}

The central finding is that exclusivity regularization has \emph{no significant effect on closed-loop cost reduction} (Cohen's $d = 0.07$, paired $t = 0.57$, $p > 0.05$; the arm without exclusivity wins 30 of 63 matched pairs) while it \emph{significantly worsens} both basin separation and state recovery. On cosine separation, the arm without exclusivity achieves a mean of 0.801 versus 0.740 with exclusivity, a large and highly significant difference (Cohen's $d = -0.97$, paired $t = -7.66$, $p < 0.001$; no-exclusivity wins 54 of 63 pairs). The separation gap is driven primarily by inter-basin cosine similarity: exclusivity raises the mean inter-basin cosine from 0.179 to 0.250, making latent representations of different basins \emph{less orthogonal}. Although the intra-basin cosine is marginally higher with exclusivity (0.990 vs.\ 0.980), this small gain in within-basin tightness is overwhelmed by the loss of between-basin discrimination.

Recovery success (M5) likewise favors the no-exclusivity arm (mean 0.549 vs.\ 0.483, $d = -0.82$, $p < 0.001$), and this advantage is consistent across all $B_{\text{proxy}}$ values and target sizes.

Exclusivity does confer a modest spectral-stability advantage: the median maximum spectral radius is 0.998 with exclusivity versus 1.002 without, and 45\% of exclusivity runs are spectrally unstable compared with 61\% without. This confirms the earlier observation that exclusivity provides an implicit spectral constraint. However, this stability advantage does not translate into any measurable downstream control benefit.

These results have a direct implication for the architecture comparison: the arrowhead arm used in the prior Stage-2 and Stage-3 decisions (\texttt{ah\_prag}) included exclusivity, which \emph{handicapped} its basin separation and recovery relative to an otherwise identical configuration without exclusivity. Therefore, the previously reported near-ties between arrowhead and block-diagonal on control metrics likely understate arrowhead's representational quality.

To test this directly, we re-ran the Stage-2 pairwise decision comparing block-diagonal (\texttt{bd\_c2}, 72 runs from the original pipeline) against arrowhead without exclusivity (\texttt{ah\_prag\_no\_excl}, 72 runs from this ablation after recovery completion). The formal decision outcome is unchanged: the pairwise cost-reduction contrast is $\Delta M_4 = -11.585$ with 95\% CI $[-34.739, +0.009]$ (includes zero), so the pre-registered rule returns no clear winner and falls back to \texttt{bd\_c2}. This is nearly identical to the original comparison with exclusivity ($\Delta M_4 = -11.593$, CI $[-34.750, +0.001]$). The bottleneck is not the arrowhead arm's quality---which is consistent across all 72 runs ($M_4 \in [0.79, 0.85]$)---but the block-diagonal arm's catastrophic outlier regime (\texttt{bd\_c2}, $\text{ts}=512$, $B_{\text{proxy}}=13$: mean $M_4 = -116.8$), which inflates bootstrap variance and prevents the confidence interval from excluding zero.

The practical conclusion is twofold. First, the arrowhead without exclusivity is strictly more reliable on closed-loop cost reduction: zero catastrophic failures, tight variance, and better basin separation and recovery. Second, the current CI-based decision rule cannot express ``arm B is more reliable'' as distinct from ``no difference,'' highlighting the need for robust effect summaries (e.g., median-based or trimmed-mean comparisons, or explicit tail-risk penalties) before the pre-registered rule can make a decisive architecture call.

\subsection{Takeaways and Current Status}

The experiments establish a coherent trajectory toward the project goal:
\begin{enumerate}
    \item \textbf{Achieved (representation identifiability):} basin-unique sparse supports at realistic basin count (13 basins) with sufficient latent capacity ($\text{ts}\geq256$), plus high threshold-free separation (intra-basin cosine $\approx 0.97$, strong intra--inter gap).
    \item \textbf{Achieved (dynamics criterion for local modeling):} long-horizon behavior is now mechanistically explained by spectral radius; SR$<1$ consistently corresponds to bounded H1000 error, while SR$>1$ yields divergence.
    \item \textbf{Achieved (local LQR readiness on benchmarks):} staged Lyapunov/Duffing evaluations show DARE feasibility rate $=1.0$ and closed-loop stability rate $=1.0$ for finalist arms, establishing that discovered local regimes admit feasible and stable local LQR controllers.
    \item \textbf{Partially achieved (architecture selection):} arrowhead vs block-diagonal remains formally unresolved under the pre-registered decision rule, even after removing the exclusivity handicap. The exclusivity ablation shows that exclusivity regularization was a net negative (significantly worse basin separation and recovery, with no control benefit). Re-running the pairwise decision with arrowhead without exclusivity yields an identical formal outcome (``no clear winner''), but this is because the decision rule's CI is dominated by catastrophic block-diagonal outliers, not because the arms are truly comparable in reliability. Arrowhead without exclusivity is strictly more reliable on $M_4$ (zero catastrophic failures, tight variance) and has better CosSep and $M_5$. The current default remains block-diagonal pending adoption of robust effect summaries (median/$M_4$ tail-risk) in the decision rule.
    \item \textbf{Open (deployment-critical):} robust label-free regime assignment and controller switching across basin transitions, where periodic re-encoding should mediate support updates and local-controller handoff.
    \item \textbf{Open (evaluation bottleneck):} current feasibility/stability metrics saturate in benchmark settings and closed-loop cost reduction can be heavy-tailed, so stronger stress tests and robust effect summaries are needed for decisive model selection.
\end{enumerate}

Overall, the narrative has shifted from ``can we learn basin-discriminative sparse representations?'' to ``can we operationalize stable, label-free controller switching across regimes?'' The first question is now strongly supported; the second is the central next milestone.
